\documentclass[10pt,a4paper]{amsart}
\usepackage[margin=19mm]{geometry}
\usepackage{amsmath,amssymb,mathtools}
\usepackage[scr=rsfs,cal=euler]{mathalfa}
\usepackage{xfrac}
\usepackage[unicode,hidelinks,bookmarks=false]{hyperref}
\newcommand{\ie}{i.e.\ }
\newcommand{\cf}{cf.\ }
\newcommand{\resp}{respectively\ }
\newcommand{\Subsection}{\subsection}
\providecommand{\nobreakdash}{\nobreak\hbox{-}}
\setlength{\emergencystretch}{2em}
\multlinegap0pt
\usepackage{amssymb}
\usepackage{bm}
\usepackage{float}
\usepackage{tikz}
\usepackage{algorithm}\usepackage{algpseudocode}
\usepackage{mathtools}
\usepackage{enumerate}
\newtheorem{corollary}{Corollary}
\newtheorem{theorem}{Theorem}
\newcommand{\balpha}{{\boldsymbol \alpha}}
\newcommand{\bx}{{\mathbf x}}
\newcommand{\bzero}{\mathbf{0}}
\title{Quantum Low-Density Parity-Check Codes from Dyadic Matrices}
\author{Anthony G\'omez-Fonseca, Gretchen L. Matthews, Kirsten D. Morris, Tefjol Pllaha}
\date{Source: arXiv:2605.01942v1 (CC BY 4.0)}
\begin{document}
\maketitle

\noindent\emph{Source and licence.} Quantum Low-Density Parity-Check Codes from Dyadic Matrices. Version arXiv:2605.01942v1 (CC BY 4.0), distributed by arXiv under the Creative Commons Attribution 4.0 International licence, http://creativecommons.org/licenses/by/4.0/, as recorded in the arXiv licence metadata for this exact version and shown on the abstract page https://arxiv.org/abs/2605.01942v1. Full licence notice: \texttt{licenses/arxiv-2605.01942-CC-BY-4.0.txt}.\par\smallskip

\noindent\textit{Editorial scope.} Counted unit: the article's theorem counting\_cycles\_thm\_6cycles\_ncxnv\_single\_dyadics (source0.tex lines 1521--1540). The article's complete original proof of it is delivered with it (source0.tex lines 1541--1566, one \texttt{proof} environment of 26 lines). No mathematics is written, repaired or re-derived: the statement and the proof are the article's own lines, in the article's own notation, and no retained line is substituted or re-flowed.  Retained uncounted as the source-local context the proof needs: lines 1079--1082; lines 1104--1107.  Nothing was omitted from any retained span, so the package carries no omission record.  No retained line contains a citation, so the wrapper carries no bibliography and no bibliography entry.  No retained line contains a figure environment, an image-inclusion command or drawing code, so no image is delivered and the package declares no asset dependency.  Editorial formatting only, in addition: an AMS \texttt{amsart} preamble that restates the article's own declarations and macros used by the retained lines, namely the environments \texttt{corollary}, \texttt{theorem} on separate counters, together with the standard packages those lines need.  The retained environments print in this document as Corollary 1, Theorem 1 in order of appearance, whereas the article prints the same items with the numbering of its own full document.  The complete native source of the article as distributed in the arXiv e-print for this exact version is bundled unchanged as \texttt{sources/199675/source0.tex}.
\begin{corollary}
\label{counting_cycles_dyadics_preimageTBCwalk_are_cycles}
Let $\tilde{G}$ be an $N$-fold graph cover of the protograph $G$ and let $W$ be a TBC walk of length $k$ in $G$. Then the inverse image of $W$ in $\tilde{G}$ is the union of $N/m$ $(km)$-cycles, where $m\geq 1$ is the order of the permutation shift of $W$ in $\mathbb{F}_2^\ell$.
\end{corollary}

\begin{equation}
\label{parity_check_matrix_single_dyadics}
D=P_{\bx_0}+P_{\bx_1}+P_{\bx_2}+\cdots+P_{\bx_{\omega-1}}.
\end{equation}

\begin{theorem}
\label{counting_cycles_thm_6cycles_ncxnv_single_dyadics}
Let $D$ be as in \eqref{parity_check_matrix_single_dyadics} and let
\begin{equation*}
A=\left\{\left(u,v,w,\bx_{u}+\bx_{v}+\bx_{w}\right)\mid u,v,w\in[\omega],u\neq v,v\neq w\right\}.
\end{equation*}
\noindent For $u,v,w,u',v',w'\in[\omega]$, let $(u,v,w,\balpha_{u,v,w})$ and $(u',v',w',\balpha_{u',v',w'}) \in A$ be such that $\balpha_{u,v,w}=\balpha_{u',v',w'}$.  
Then this repetition $\balpha_{u,v,w}=\balpha_{u',v',w'}$ lifts to a collection of 6-cycles in the Tanner graph if $u\neq u'$ and $w\neq w'$. The total number of 6-cycles in the Tanner graph, $\mathcal{N}_6$, is given by
\begin{equation}
\label{counting_cycles_formula_num6cycles_ncxnv_single_dyadics}
\mathcal{N}_6=\frac{2^\ell}{6}\cdot\left(\mathcal{R}_{A}^{*c_1}+\mathcal{R}_{A}^{*c_2}+\mathcal{R}_{A}^{*c_3}+\mathcal{R}_{A}^{*c_4}+\mathcal{R}_{A}^{*nc}\right),
\end{equation}
where $\mathcal{R}_{A}^{*c_1}$, $\mathcal{R}_{A}^{*c_2}$, $\mathcal{R}_{A}^{*c_3}$, $\mathcal{R}_{A}^{*c_4}$, and $\mathcal{R}_{A}^{*nc}$ are the numbers of repetitions $\balpha_{u,v,w}=\balpha_{u',v',w'}$ in the mset $A$ satisfying Conditions 1, 2, 3, 4, and not satisfying either, respectively, and where
\begin{itemize}
\item Condition 1: $u=w'$, $v=v'$, and $w=u'$;
\item Condition 2: $v=u'$ and $w=v'$;
\item Condition 3: $u=w$ and $w'=u'$; and
\item Condition 4: $u=v'$ and $v=w'$.
\end{itemize} 
\end{theorem}

\begin{proof}
For $u,v,w,u',v',w'\in[\omega]$, let $(u,v,w,\balpha_{u,v,w})$ and $(u',v',w',\balpha_{u',v',w'})$ be two elements in $A$ such that $\balpha_{u,v,w}=\balpha_{u',v',w'}$. The walk of length 6 corresponding to this repetition is given by $W=[\bx_{u},\bx_{v},\bx_{w},\bx_{w'},\bx_{v'},\bx_{u'}]$. If $u\neq u'$ and $w\neq w'$, then this walk is actually a TBC walk with permutation shift $\balpha_{u,v,w}+\balpha_{u',v',w'}=\bzero$, so by Corollary \ref{counting_cycles_dyadics_preimageTBCwalk_are_cycles}, its inverse image is the union of $N/m$ 6-cycles, where $m\geq 1$ is the order of the permutation shift of $W$ in $\mathbb{F}_2^\ell$. Notice that $m\in\{1,2\}$, but since $m=2$ is not possible by the argument before the statement of Theorem \ref{counting_cycles_thm_6cycles_ncxnv_single_dyadics}, we only study the case $m=1$. 

Let us focus on the indices $(u,v,w,w',v',u')$. When computing all the equivalent walks to $W$, there are five sets of conditions that arise:
\begin{itemize}
\item Condition 1: $u=w'$, $v=v'$, and $w=u'$;
\item Condition 2: $v=u'$ and $w=v'$;
\item Condition 3: $u=w$ and $w'=u'$;
\item Condition 4: $u=v'$ and $v=w'$; and
\item Condition 5: $u=w=v'$ and $v=w'=u'$.
\end{itemize} 
If Condition 1 is assumed, the tuple of indices $(u,v,w,w',v',u')$ becomes $(u,v,w,u,v,w)$. Since the walk $[\bx_{u},\bx_{v},\bx_{w}]$ is not closed, the walk corresponding to the tuple $(u,v,w,u,v,w)$, $W$, is not the double traversal of a TBC walk of length 3 (which would require $m=2$), so the tuple actually contributes to the number of 6-cycles in the Tanner graph. When computing all the walks equivalent to $W$, we obtain exactly 6 distinct elements $(u,v,w,u,v,w)$, $(v,w,u,v,w,u)$, $(w,u,v,w,u,v)$, $(u,w,v,u,w,v)$, $(v,u,w,v,u,w)$, and $(w,v,u,w,v,u)$, which means that the repetition $\balpha_{u,v,w}=\balpha_{u',v',w'}$ is counted six times in the construction of $A$. The number of 6-cycles arising from this repetition is $2^\ell$.

If Condition 2 is assumed, the tuple of indices $(u,v,w,w',v',u')$ becomes $(u,v,w,w',w,v)$.  
When computing all the walks equivalent to $W$, we obtain 6 distinct elements $(u,v,w,w',w,v)$, $(v,w,w',w,v,u)$, $(w,w',w,v,u,v)$, $(w',w,v,u,v,w)$, $(w,v,u,v,w,w')$, and $(v,u,v,w,w',w)$, which means that the repetition $\balpha_{u,v,w}=\balpha_{u',v',w'}$ is counted six times in the construction of $A$. The number of 6-cycles arising from this repetition is $2^\ell$.

If Condition 3 is assumed, the tuple of indices $(u,v,w,w',v',u')$ becomes $(u,v,u,u',v',u')$. 
When computing all the walks equivalent to $W$, we obtain 6 distinct elements $(u,v,u,u',v',u')$, $(v,u,u',v',u',u)$, $(u,u',v',u',u,v)$, $(u',v',u',u,v,u)$, $(v',u',u,v,u,u')$, and $(u',u,v,u,u',v')$, which means that the repetition $\balpha_{u,v,w}=\balpha_{u',v',w'}$ is counted six times in the construction of $A$. The number of 6-cycles arising from this repetition is $2^\ell$.

If Condition 4 is assumed, the tuple of indices $(u,v,w,w',v',u')$ becomes $(u,v,w,v,u,u')$. 
When computing all the walks equivalent to $W$, we obtain 6 distinct elements $(u,v,w,v,u,u')$, $(v,w,v,u,u',u)$, $(w,v,u,u',u,v)$, $(v,u,u',u,v,w)$, $(u,u',u,v,w,v)$, and $(u',u,v,w,v,u)$, which means that the repetition $\balpha_{u,v,w}=\balpha_{u',v',w'}$ is counted six times in the construction of $A$. The number of 6-cycles arising from this repetition is $2^\ell$.

If Condition 5 is assumed, the tuple of indices $(u,v,w,w',v',u')$ becomes $(u,v,u,v,u,v)$. The corresponding TBC walk would be formed by traversing a TBC walk of length 2 three times, but we have argued why this is not possible. Finally, if none of these conditions are satisfied, then we also obtain 6 distinct elements $(u,v,w,w',v',u')$, $(v,w,w',v',u',u)$, $(w,w',v',u',u,v)$, $(w',v',u',u,v,w)$, $(v',u',u,v,w,w')$, and $(u',u,v,w,w',v')$, when computing all the walks equivalent to $W$. This means that the repetition $\balpha_{u,v,w}=\balpha_{u',v',w'}$ is counted six times in the construction of $A$. The number of 6-cycles arising from this repetition is $2^\ell$.

Therefore, if $\mathcal{R}_{A}^{*c_1}$, $\mathcal{R}_{A}^{*c_2}$, $\mathcal{R}_{A}^{*c_3}$, $\mathcal{R}_{A}^{*c_4}$, and $\mathcal{R}_{A}^{*nc}$ are the numbers of repetitions $\balpha_{u,v,w}=\balpha_{u',v',w'}$ in the mset $A$ satisfying Conditions 1, 2, 3, 4, and not satisfying either, respectively, then \eqref{counting_cycles_formula_num6cycles_ncxnv_single_dyadics} follows and we conclude the proof.
\end{proof}

\end{document}
